\chapter{Modality Selection and Appropriateness Quick Reference}
\label{app:modality_selection}

Selecting the right imaging examination depends on the clinical question, urgency, patient age and pregnancy status, renal function, implanted devices, and local availability. This appendix summarizes the comparative strengths, limitations, and primary clinical roles of the major imaging modalities. Formal appropriateness decisions should consult current evidence-based referral guidance such as the American College of Radiology (ACR) \textit{Appropriateness Criteria}.

\section{Cross-Modality Comparison}

\begin{longtable}{@{}p{0.16\textwidth}p{0.14\textwidth}p{0.32\textwidth}p{0.30\textwidth}@{}}
\caption{Comparison of core diagnostic imaging modalities.}
\label{tab:app_modality_comparison} \\
\toprule
\textbf{Modality} & \textbf{Ionizing Radiation} & \textbf{Primary Strengths} & \textbf{Key Limitations} \\
\midrule
\endfirsthead
\toprule
\textbf{Modality} & \textbf{Ionizing Radiation} & \textbf{Primary Strengths} & \textbf{Key Limitations} \\
\midrule
\endhead
Radiography (X-ray) & Yes (very low) & Rapid, portable, low cost; bone alignment, pneumothorax, lines/tubes, bowel obstruction & 2D projection overlap; limited soft-tissue contrast \\
Ultrasound (US) & No & Real-time, portable, no radiation; fluid vs.\ solid, biliary/renal, obstetrics, vascular Doppler & Operator-dependent; limited by gas, bone, and deep habitus \\
Computed Tomography (CT) & Yes (moderate) & Fast isotropic 3D anatomy; trauma, acute hemorrhage, lung parenchyma, CTA, stone/staging & Ionizing radiation; lower soft-tissue contrast than MRI \\
Magnetic Resonance Imaging (MRI) & No & Superior soft-tissue contrast; brain, spine, joints, liver, pelvis, multiparametric oncology & Longer scan time, motion sensitivity, implant/claustrophobia constraints \\
Nuclear Medicine / SPECT & Yes (low--mod) & Functional/physiological perfusion, bone turnover, thyroid, hepatobiliary, V/Q & Lower spatial resolution; radiopharmaceutical availability \\
PET / CT (FDG) & Yes (mod--high) & Whole-body metabolic staging, treatment response, occult infection/inflammation & Physiological uptake mimics; limited resolution for sub-cm lesions \\
\bottomrule
\end{longtable}

\section{Common Clinical Presentations}

\begin{longtable}{@{}p{0.26\textwidth}p{0.26\textwidth}p{0.40\textwidth}@{}}
\caption{Typical first-line imaging choices by clinical scenario.}
\label{tab:app_clinical_selection} \\
\toprule
\textbf{Clinical Scenario} & \textbf{Typical Initial Modality} & \textbf{Key Follow-Up / Alternative} \\
\midrule
\endfirsthead
\toprule
\textbf{Clinical Scenario} & \textbf{Typical Initial Modality} & \textbf{Key Follow-Up / Alternative} \\
\midrule
\endhead
Acute focal neurological deficit & Non-contrast head CT (NCCT) + CTA & Brain MRI (DWI/FLAIR) if NCCT negative or wake-up stroke \\
Suspected pulmonary embolism & CT pulmonary angiography (CTPA) & V/Q scan in severe renal failure or selected pregnancy cases \\
Suspected aortic dissection & CTA chest/abdomen/pelvis & Transesophageal echo or MRA if CTA contraindicated \\
Right upper quadrant pain / biliary & Abdominal ultrasound & MRCP or HIDA scan if ultrasound equivocal \\
Renal colic / ureteral stone & Non-contrast low-dose CT KUB & Renal/bladder ultrasound in pregnancy or pediatrics \\
Suspected appendicitis (adult) & Contrast-enhanced CT abdomen/pelvis & Ultrasound or MRI in pregnancy and young children \\
Suspected cauda equina / cord compression & Urgent spine MRI & CT myelography if MRI contraindicated \\
Occult hip fracture (negative X-ray) & Hip MRI without contrast & CT if MRI unavailable or contraindicated \\
\bottomrule
\end{longtable}
